\chapter{Birch and Swinnerton-Dyer}\label{ch:bsd}

\begin{center}\itshape
A curve with many rational points should have many points modulo $p$, for most $p$.\\
\upshape(the intuition behind the conjecture, in our words)
\end{center}

In the early 1960s, in the Mathematical Laboratory at Cambridge, Bryan Birch and Peter Swinnerton-Dyer kept the EDSAC~2 computer busy with a question that
looks almost naive. Take an elliptic curve over $\Q$, say $y^2=x^3-Dx$. For each prime $p$, count its points modulo $p$: call the number $N_p$. Hasse's theorem
says that $N_p$ is close to $p+1$, within $2\sqrt p$. Now multiply the ratios $N_p/p$ over all primes up to $x$. Does this product know anything about the
rational points of the curve?

It did. For curves with only finitely many rational points the product seemed to stay bounded; for curves of rank $1$ it grew like $\log x$; for rank $2$ like
$(\log x)^2$. Birch and Swinnerton-Dyer turned the pattern into a conjecture about the $L$-function of the curve at the point $s=1$, refined it into an exact
formula, and published it in two papers in Crelle's journal (1963 and 1965). In 2000 the Clay Mathematics Institute put it on its list of seven Millennium
Problems. Sixty years after the first computations, it is proved for the curves whose $L$-function vanishes to order at most $1$ at $s=1$, and for no curve of rank $2$ or more
does anyone know a general reason why the analytic and algebraic ranks should agree.

Tome~12 built the stage: the group law and Mordell's theorem (Chapter~12), Hasse's bound (Chapter~13), and modular forms and the modularity theorem
(Chapter~14). It also stated the conjecture in a remark and proved, with the congruent number curves $E_n:y^2=x^3-n^2x$, that $1$ is not a congruent number.
This chapter does what Tome~12 did not: it counts the points of the curves $E_n$ exactly, by Jacobi sums, as in Birch and Swinnerton-Dyer's first paper; it
proves that the Euler product converges where it should; it shows that exactly half of the $a_p$ of these curves vanish; and it runs the experiment again,
on four curves of ranks $0$ to $3$, up to $2\cdot10^5$.

\begin{think}
Count the points of $y^2=x^3-x$ modulo $5$ and modulo $13$, including the point at infinity. Both counts are divisible by $8$. Is that an accident? Before
reading on, try $p=17$ and $p=29$.
\end{think}

\section{The conjecture and the evidence}

Let $E$ be an elliptic curve over $\Q$, given by a Weierstrass equation with integer coefficients and minimal discriminant $\Delta$. For a prime $p\nmid\Delta$
(a \emph{good} prime) the reduction of $E$ modulo $p$ is an elliptic curve over $\mathbb F_p$, with $N_p$ points (counting $O$), and we put $a_p=p+1-N_p$. By
Hasse's theorem (Tome~12, Corollary~13.14), $|a_p|\le2\sqrt p$. For a bad prime we put $a_p=p-(\text{number of affine solutions modulo }p)$, a convention
which gives $a_p\in\{0,\pm1\}$ for a minimal equation; we only need that there are finitely many bad primes. The \emph{$L$-function} of $E$ is
\[
L(E,s)=\prod_{p\nmid\Delta}\frac1{1-a_pp^{-s}+p^{1-2s}}\;\prod_{p\mid\Delta}\frac1{1-a_pp^{-s}} .
\]
Formally, at $s=1$ the factor at a good prime is $\bigl(1-\frac{a_p}p+\frac1p\bigr)^{-1}=\frac p{N_p}$. So if $E$ has many points modulo $p$ for most $p$, the
factors are small and $L(E,1)$ ``wants'' to vanish.

\begin{center}
\emph{The Birch and Swinnerton-Dyer conjecture (rank part).} $L(E,s)$ has a zero at $s=1$ of order exactly the rank $r$ of $E(\Q)$.
\end{center}

The product does not converge at $s=1$, so the statement needs $L(E,s)$ to exist there. That it does is a deep theorem: by the modularity theorem
(Tome~12, Chapter~14), $L(E,s)$ is the $L$-function of a modular form of weight $2$ and level $N$, the \emph{conductor} of $E$, so it extends to an entire
function with a functional equation relating $s$ and $2-s$. Before 1995 the conjecture could only be stated for special curves, such as the curves with
complex multiplication that this chapter studies, for which Deuring had proved the continuation in the 1950s. The full conjecture also predicts the leading
coefficient: if $L(E,s)\sim C(s-1)^r$, then
\[
C=\frac{\Omega_E\,R_E\,|\operatorname{Sha}(E)|\prod_pc_p}{|E(\Q)_{\mathrm{tors}}|^2},
\]
where $\Omega_E$ is the real period (an integral of $\frac{dx}{2y+a_1x+a_3}$ over $E(\R)$), $R_E$ the regulator (a determinant of canonical heights of a basis
of the free part, so $R_E=1$ when $r=0$), $c_p$ the Tamagawa numbers (small integers attached to the bad primes), and $\operatorname{Sha}(E)$ the Tate--Shafarevich group (written with the Cyrillic letter sha in most texts),
conjectured to be finite. For a curve of rank $0$ the formula says that $L(E,1)/\Omega_E$ is a rational number with a predicted arithmetic meaning; the
notebook computes $L(E,1)$ from its rapidly converging series and finds, for the curve $y^2+y=x^3-x^2$ of Tome~12,
\[
L(E,1)=0.253841\ldots,\qquad\Omega_E=6.346046\ldots,\qquad\frac{L(E,1)}{\Omega_E}=0.0400000\ldots=\frac1{25},
\]
which is $\frac1{|E(\Q)_{\mathrm{tors}}|^2}$ for its torsion group $\Z/5$, with all the other factors equal to $1$. For $y^2=x^3-x$ it finds
$L(E,1)/\Omega_E=0.125000\ldots=\frac18$; with $|E(\Q)_{\mathrm{tors}}|=4$ the formula then asks that $|\operatorname{Sha}(E)|\prod c_p=2$.

\paragraph{The evidence.} Figure~\ref{fig:bsd} redoes the original experiment with four curves taken from the tables, of ranks $0$, $1$, $2$ and $3$: our
$y^2=x^3-x$ (rank $0$, by Tome~12, Theorem~12.15), and the curves of smallest conductor of ranks $1$, $2$ and $3$, of conductors $37$, $389$ and $5077$. The
four curves are visibly sorted by rank. Fitting straight lines against $\log\log x$ over the primes $10^3<p<2\cdot10^5$ gives slopes $-0.00$, $1.03$,
$2.17$ and $2.91$; over $10^4<p<2\cdot10^5$ the same fits give $0.10$, $1.01$, $2.35$ and $2.66$. The slopes are right to within a few tenths, and they move
by a few tenths when the window moves: $\log\log x$ grows so slowly that the whole computation covers less than one unit of it. Birch and Swinnerton-Dyer
saw a pattern of this kind, not a limit, and nobody today can compute far enough to see a limit.

\begin{figure}[ht]
\centering
\fig{ch10-bsd}
\caption{$\log\prod_{p\le x}N_p/p$ over the good primes $p\le x$, for $x$ up to $2\cdot10^5$, plotted against $\log\log x$, for four curves of ranks $0$ to $3$
(conductors $32$, $37$, $389$, $5077$). The values at $x=2\cdot10^5$ are $0.905$, $4.593$, $6.766$ and $9.047$. If the product behaves like $C(\log x)^r$,
the curves become straight lines of slope $r$. Point counts computed for this book.}
\label{fig:bsd}
\end{figure}

\paragraph{The heuristic.} Why $(\log x)^r$? Two reasons, one arithmetic and one analytic. The arithmetic one is the epigraph: a rational point of infinite
order reduces to a point modulo almost every $p$, and its multiples give many points there, so a large rank should push $N_p$ above $p+1$ on average. The
analytic one is Mertens' theorem. The Riemann zeta function has a pole of order $1$ at $s=1$, and its partial Euler products satisfy
$\prod_{p\le x}(1-\frac1p)^{-1}\sim e^\gamma\log x$ (Tome~5, Theorem~1.6). If the partial Euler products of $L(E,s)$ at $s=1$ behaved in the same way, a zero
of order $r$ would give $\prod_{p\le x}\frac p{N_p}\approx\frac C{(\log x)^r}$. This is only an analogy: for $L(E,s)$ the line $\operatorname{Re}s=1$ is the centre of the critical strip
$\frac12<\operatorname{Re}s<\frac32$, half a unit away from the half-plane where the Euler product converges (Proposition~\ref{prop:eulerconv}), and
Remark~\ref{rem:goldfeld} explains what is known.

\section{Counting points with Jacobi sums}

For a prime $p$ and characters $\chi,\psi$ of $\mathbb F_p^\times$, extended by $\chi(0)=0$ when $\chi$ is non-trivial, the \emph{Jacobi sum} is
$J(\chi,\psi)=\sum_{u\in\mathbb F_p}\chi(u)\psi(1-u)$, and the \emph{Gauss sum} is $g(\chi)=\sum_{u\in\mathbb F_p}\chi(u)\zeta^u$ with $\zeta=e^{2\pi i/p}$.
For a non-trivial $\chi$, $|g(\chi)|=\sqrt p$ (Tome~5, Theorem~3.5: every non-trivial character modulo a prime is primitive). We write $\rho$ for the
Legendre symbol, the character of order $2$.

\begin{lemma}\label{lem:jacabs}
If $\chi$, $\psi$ and $\chi\psi$ are non-trivial characters of $\mathbb F_p^\times$, then $g(\chi)g(\psi)=J(\chi,\psi)\,g(\chi\psi)$, and $|J(\chi,\psi)|=\sqrt p$.
\end{lemma}
\begin{proof}
Expand and group by $u=s+t$:
\[
g(\chi)g(\psi)=\sum_{s,t}\chi(s)\psi(t)\zeta^{s+t}=\sum_u\zeta^u\sum_{s+t=u}\chi(s)\psi(t).
\]
For $u\ne0$, substitute $s=us'$, $t=ut'$ with $s'+t'=1$: the inner sum is $\chi\psi(u)J(\chi,\psi)$. For $u=0$ it is
$\sum_s\chi(s)\psi(-s)=\psi(-1)\sum_s\chi\psi(s)=0$, since $\chi\psi$ is non-trivial. Hence $g(\chi)g(\psi)=J(\chi,\psi)\sum_{u}\chi\psi(u)\zeta^u=J(\chi,\psi)g(\chi\psi)$.
Taking absolute values, $p=|J(\chi,\psi)|\sqrt p$.
\end{proof}

Now let $p\equiv1\pmod4$. The group $\mathbb F_p^\times$ is cyclic of order divisible by $4$, so it has a character $\chi$ of order $4$ (send a generator to $i$).
Then $\chi^2=\rho$, $\chi^3=\overline\chi$, and $\chi$ takes values in $\{0,\pm1,\pm i\}$, so every Jacobi sum built from $\chi$ and $\rho$ lies in $\Z[i]$.

\begin{lemma}\label{lem:quarticap}
Let $p\equiv1\pmod4$ and $\chi$ a character of order $4$ of $\mathbb F_p^\times$. For $E_1:y^2=x^3-x$, $a_p=-2\operatorname{Re}J(\chi,\rho)$. Writing
$J(\chi,\rho)=A+Bi$ with $A,B\in\Z$, we have $A^2+B^2=p$ and $a_p=-2A$.
\end{lemma}
\begin{proof}
For each $x$, the equation $y^2=c$ has $1+\rho(c)$ solutions $y$ (this is $1$ for $c=0$). So the number of affine points is $\sum_x(1+\rho(x^3-x))$, and with
the point at infinity, $a_p=p+1-N_p=-\sum_x\rho(x^3-x)=-\sum_x\rho(x)\rho(x^2-1)$.

Group the non-zero $x$ by $u=x^2$. Each non-zero square $u$ comes from two values $\pm x$, and $\rho(-x)=\rho(x)$ because $\rho(-1)=1$ for $p\equiv1\pmod4$.
Moreover $\rho(x)=\chi(x)^2=\chi(x^2)=\chi(u)$. Since $1+\rho(u)$ is $2$ on the non-zero squares and $0$ on the non-squares,
\[
\sum_x\rho(x)\rho(x^2-1)=\sum_{u\ne0}\bigl(1+\rho(u)\bigr)\chi(u)\rho(u-1)=\sum_u\chi(u)\rho(1-u)+\sum_u\overline\chi(u)\rho(1-u),
\]
using $\rho(u-1)=\rho(-1)\rho(1-u)=\rho(1-u)$ and $\rho\chi=\chi^3=\overline\chi$. The two sums are $J(\chi,\rho)$ and $J(\overline\chi,\rho)=\overline{J(\chi,\rho)}$,
so $a_p=-2\operatorname{Re}J(\chi,\rho)=-2A$. Finally $\chi$, $\rho$ and $\chi\rho=\overline\chi$ are non-trivial, so $A^2+B^2=|J(\chi,\rho)|^2=p$ by
Lemma~\ref{lem:jacabs}.
\end{proof}

The lemma determines $|a_p|$: $p=A^2+B^2$, and $p$ has only one such decomposition up to order and signs ($p$ has exactly $8$ representations as a sum of two
squares, Tome~1, Theorem~10.5). But which of the two squares is $A^2$, and what is the sign of $A$? The group structure answers both.

\begin{lemma}\label{lem:eightdiv}
If $p\equiv1\pmod4$, then $8$ divides $N_p$, the number of points of $y^2=x^3-x$ over $\mathbb F_p$.
\end{lemma}
\begin{proof}
Let $i\in\mathbb F_p$ with $i^2=-1$, and $P=(i,1-i)$. It lies on the curve: $(1-i)^2=1-2i+i^2=-2i$ and $i^3-i=-i-i=-2i$. The tangent at $P$ has slope
$\lambda=\frac{3x^2-1}{2y}=\frac{-4}{2(1-i)}=\frac{-2}{1-i}=-(1+i)$, so the $x$-coordinate of $2P$ is $\lambda^2-2x=(1+i)^2-2i=0$, and $2P=(0,0)$
(the only point with $x=0$), which has order $2$. So $P$ has order $4$ and $\langle P\rangle=\{O,P,(0,0),-P\}$. The point $(1,0)$ has order $2$ and is not
in $\langle P\rangle$, since $\pm P$ have $x=i\ne1$. So the subgroup generated by $P$ and $(1,0)$ is $\langle P\rangle\cup\bigl((1,0)+\langle P\rangle\bigr)$,
of order $8$, and Lagrange's theorem gives $8\mid N_p$ (the points form a group by Tome~12, Theorem~12.2).
\end{proof}

\begin{theorem}\label{thm:countEn}
Let $n\ge1$ and $p\nmid2n$ a prime, and let $N_p$ be the number of points of $E_n:y^2=x^3-n^2x$ over $\mathbb F_p$, $a_p=p+1-N_p$.
\begin{enumerate}[label=(\alph*)]
\item If $p\equiv3\pmod4$, then $a_p=0$.
\item If $p\equiv1\pmod4$, write $p=a^2+b^2$ with $a$ odd and $a\equiv\frac{p+1}2\pmod4$; this determines $a$ uniquely. Then
\[
a_p=2\Bigl(\frac np\Bigr)a .
\]
\end{enumerate}
In particular $a_p\ne0$ in case (b).
\end{theorem}
\begin{proof}
\emph{Reduction to $n=1$.} The substitution $x=nu$ is a bijection of $\mathbb F_p$, and $x^3-n^2x=n^3(u^3-u)$, so
$a_p(E_n)=-\sum_x\rho(x^3-n^2x)=-\rho(n)^3\sum_u\rho(u^3-u)=\rho(n)\,a_p(E_1)$. (Case (a) for $E_n$ is also in the proof of Tome~12, Proposition~12.14.)

(a) If $p\equiv3\pmod4$, then $\rho(-1)=-1$, and $u\mapsto-u$ changes the sign of $\rho(u^3-u)$, so the sum is $0$.

(b) By Lemma~\ref{lem:quarticap}, $a_p(E_1)=-2A$ with $A^2+B^2=p$, so $N_p=p+1+2A$, and by Lemma~\ref{lem:eightdiv}, $2A\equiv-(p+1)\pmod8$, that is
$A\equiv-\frac{p+1}2\pmod4$. Since $p\equiv1\pmod4$, $\frac{p+1}2$ is odd, so $A$ is odd. Put $a=-A$: then $a$ is odd, $a\equiv\frac{p+1}2\pmod4$, $a^2+B^2=p$,
and $a_p(E_1)=2a$. For uniqueness: the representations of $p$ as a sum of two squares are the $8$ pairs $(\pm a_0,\pm b_0)$, $(\pm b_0,\pm a_0)$ with $a_0$ odd
and $b_0$ even (Tome~1, Theorem~10.5, and $p$ odd), so the odd member is $\pm a_0$, and exactly one of the two signs is $\equiv\frac{p+1}2\pmod4$. Finally
$a$ is odd, so $a_p\ne0$.
\end{proof}

For $p=5$: $5=1+4$, $\frac{p+1}2=3$, so $a=-1$ and $N_5=5+1+2=8$. For $p=13$: $13=9+4$, $\frac{p+1}2=7\equiv3$, so $a=3$, $a_{13}=6$ and $N_{13}=8$. For
$p=17$: $a=1$, $N_{17}=16$. The notebook checks the theorem for every prime below $2\cdot10^5$. This is the kind of count that Birch and Swinnerton-Dyer
needed in their first paper, which treats the curves $y^2=x^3-Dx$; and it has a famous ancestor. The last entry of Gauss's mathematical diary, dated
9~July~1814, records a count of the solutions of $x^2+y^2+x^2y^2\equiv1\pmod p$ in terms of the decomposition $p=a^2+b^2$, for a curve closely related
to the lemniscate; the method above, Jacobi sums with values in $\Z[i]$, is the modern form of his.

\section{What is proved}

\paragraph{Convergence.} The product defining $L(E,s)$ converges in the half-plane where Hasse's bound makes it converge, and no further.

\begin{proposition}\label{prop:eulerconv}
Let $E$ be an elliptic curve over $\Q$. Define $a_n$ for $n\ge1$ by $a_1=1$, $a_{mn}=a_ma_n$ for coprime $m,n$, and for each prime $p$
\[
a_{p^{k+1}}=a_pa_{p^k}-\varepsilon_p\,p\,a_{p^{k-1}}\quad(k\ge1),\qquad\varepsilon_p=\begin{cases}1&p\text{ good},\\0&p\text{ bad}.\end{cases}
\]
Then $|a_n|\le d(n)\sqrt n$ for every $n$ coprime to the bad primes, and for $\operatorname{Re}s>\frac32$ the series $\sum a_nn^{-s}$ converges absolutely and
\[
\sum_{n\ge1}\frac{a_n}{n^s}=\prod_{p\text{ good}}\frac1{1-a_pp^{-s}+p^{1-2s}}\prod_{p\text{ bad}}\frac1{1-a_pp^{-s}}=L(E,s),
\]
a holomorphic function without zeros in that half-plane.
\end{proposition}
\begin{proof}
Let $p$ be good. By Hasse's bound $a_p^2\le4p$, so the polynomial $X^2-a_pX+p$ has roots $\alpha$, $\overline\alpha$ that are complex conjugate (or equal and
real), with $\alpha\overline\alpha=p$, hence $|\alpha|=\sqrt p$. Put $b_k=\sum_{j=0}^k\alpha^j\overline\alpha^{k-j}$. Then $b_0=1$, $b_1=\alpha+\overline\alpha=a_p$,
and $b_{k+1}=(\alpha+\overline\alpha)b_k-\alpha\overline\alpha\,b_{k-1}$ (multiply out), the same recursion as $a_{p^k}$; so $a_{p^k}=b_k$ and $|a_{p^k}|\le(k+1)p^{k/2}$.
For coprime $m,n$ coprime to the bad primes, multiplicativity and $d(mn)=d(m)d(n)$ give $|a_n|\le d(n)\sqrt n$. At a bad prime, $a_{p^k}=a_p^k$ and
$|a_p|\le p$ (there are at most $2p$ affine points), so $|a_{p^k}|\le p^k$; let $B$ be the product of the bad primes.

Write $n=n_1n_2$ with $n_2$ composed of bad primes and $n_1$ coprime to $B$. For $\sigma=\operatorname{Re}s>\frac32$,
\[
\sum_n\frac{|a_n|}{n^\sigma}\le\sum_{n_1}\frac{d(n_1)}{n_1^{\sigma-1/2}}\prod_{p\mid B}\sum_{k\ge0}p^{k(1-\sigma)}<\infty,
\]
since $\sum d(n)n^{-t}=\zeta(t)^2<\infty$ for $t>1$ and $p^{1-\sigma}<1$. Tome~4, Theorem~10.3, applied to the multiplicative function $a_n$, gives the Euler
product with factors $\sum_ka_{p^k}p^{-ks}$. At a good prime, the generating function of the $b_k$ is
$\sum_kb_kT^k=\frac1{(1-\alpha T)(1-\overline\alpha T)}=\frac1{1-a_pT+pT^2}$ for $|T|<p^{-1/2}$, and at a bad prime it is $\frac1{1-a_pT}$; put $T=p^{-s}$,
with $|T|=p^{-\sigma}<p^{-3/2}$. A uniformly convergent Dirichlet series is holomorphic, and the product has no zero: each factor is the inverse of a
non-zero number, and $\sum_p|\log(1-\alpha p^{-s})|<\infty$ because $|\alpha p^{-s}|\le p^{1/2-\sigma}$ and $\sum_pp^{1/2-\sigma}<\infty$.
\end{proof}

For the modular form of Tome~12, Proposition~14.11, the coefficients $c_n$ of $\eta(\tau)^2\eta(11\tau)^2$ are exactly these $a_n$ for $y^2+y=x^3-x^2$
(Exercise~\ref{ex:anbound}). The point $s=1$ is a distance $\frac12$ to the left of the region of Proposition~\ref{prop:eulerconv}; the functional equation
of Tome~12, Theorem~14.9, sends $s=1$ to itself, so it is the centre of symmetry of $L(E,s)$, exactly as $s=\frac12$ is for $\zeta$.

\paragraph{Half the $a_p$ vanish.} For a curve without complex multiplication, $a_p=0$ is rare: the primes with $a_p=0$ have density $0$ (a theorem of Serre,
1981). For our curves, Theorem~\ref{thm:countEn} gives a sharp and opposite answer.

\begin{corollary}\label{cor:cmhalf}
For $E_n$, the good primes with $a_p=0$ are exactly the primes $p\equiv3\pmod4$ with $p\nmid n$. Among the primes up to $x$, the proportion of those with
$a_p=0$ tends to $\frac12$.
\end{corollary}
\begin{proof}
The first statement is Theorem~\ref{thm:countEn}. The second is the prime number theorem for arithmetic progressions (Tome~5, Theorem~5.8): the primes
$\equiv3\pmod4$ up to $x$ number $\sim\frac12\frac x{\log x}$, and the finitely many primes dividing $2n$ do not matter.
\end{proof}

What about the other half? For a split prime $p\equiv1\pmod4$, write $a_p=2\sqrt p\cos\theta_p$ with $\theta_p\in[0,\pi]$. By Theorem~\ref{thm:countEn},
$\cos\theta_p=\pm\frac a{\sqrt p}$, where $a+bi$ is a Gaussian prime of norm $p$: the angle $\theta_p$ is, up to symmetries, the argument of a Gaussian prime.
Figure~\ref{fig:cm} shows the $8977$ split primes below $2\cdot10^5$. The histogram of $\cos\theta_p$ is not Tome~12's semicircle but the \emph{arcsine}
density $\frac1{\pi\sqrt{1-t^2}}$, which is what one gets when $\theta_p$ is uniform on $[0,\pi]$.

\begin{figure}[ht]
\centering
\fig{ch10-cm}
\caption{Histogram ($20$ bins) of $a_p/(2\sqrt p)$ for $y^2=x^3-x$ over the $8977$ primes $p\equiv1\pmod4$ below $2\cdot10^5$, against the arcsine density
(solid) and the Sato--Tate density of curves without complex multiplication (dashed). The two outer bins have density $1.43$, where the arcsine density
is infinite at $\pm1$. Point counts computed for this book.}
\label{fig:cm}
\end{figure}

\begin{remark}[Hecke, and the CM Sato--Tate law]\label{rem:hecke}
Hecke proved in 1918--1920, with $L$-functions of characters of the ideal classes he called \emph{Gr\"o\ss encharaktere}, that the arguments of the Gaussian
primes are equidistributed: for $0\le\alpha<\beta\le\frac\pi2$, the primes $\pi=a+bi$ with $a>0$, $b\ge0$ and $\arg\pi\in[\alpha,\beta]$ have density
$\frac{\beta-\alpha}{\pi/2}$ among all such primes. In 1953 Deuring showed that the $L$-function of a curve with complex multiplication is such a Hecke
$L$-function. Combining the two, the angles $\theta_p$ of $E_n$ are equidistributed in $[0,\pi]$ along the split primes, so that over all primes their law is
$\frac12\delta_{\pi/2}+\frac12\cdot\frac{d\theta}\pi$: half of the mass sits at $\frac\pi2$ (the $a_p=0$ of Corollary~\ref{cor:cmhalf}) and half is uniform.
Compare Tome~12's remark on Sato--Tate, where the law is $\frac2\pi\sin^2\theta\,d\theta$ for every curve without complex multiplication.
Exercise~\ref{ex:cmmoments} shows how the moments tell the two laws apart.
\end{remark}

\begin{remark}[what is proved about the conjecture]\label{rem:bsdproved}
Coates and Wiles (1977) proved that if $E$ has complex multiplication and $L(E,1)\ne0$, then $E(\Q)$ is finite; our $E_n$ are such curves. Gross and Zagier
(1986) related $L'(E,1)$ to the height of a Heegner point, and Kolyvagin (1988--1989) used Heegner points to bound Selmer groups. Together with the modularity
theorem they give: \emph{if $L(E,s)$ vanishes to order $0$ or $1$ at $s=1$, then the rank of $E(\Q)$ is that order and $\operatorname{Sha}(E)$ is finite}. For a specific curve
one can sometimes go further by computing both sides: Buhler, Gross and Zagier (1985) proved that $L(E,s)$ vanishes to order exactly $3$ for the curve of
conductor $5077$ in Figure~\ref{fig:bsd}, which has rank $3$. But no theorem says that a zero of order $2$ or more forces even one rational point of
infinite order; and the order of $\operatorname{Sha}(E)$ is not known to be finite for a single curve of rank $2$ or more.
\end{remark}

\begin{remark}[on average]\label{rem:averages}
Order the curves $y^2=x^3+Ax+B$ by $\max(4|A|^3,27B^2)$. Bhargava and Shankar (2013, from the average size of $5$-Selmer groups) proved that the average rank
is less than $0.885$, that at least $20.62\%$ of the curves have rank $0$, and at least $83.75\%$ have rank $0$ or $1$; with Skinner and Zhang (2014), Bhargava
showed that at least $66.48\%$ of curves satisfy the rank part of the conjecture and have finite $\operatorname{Sha}(E)$. The minimalist conjecture is that half the curves have
rank $0$, half rank $1$, and $0\%$ rank $2$ or more.
\end{remark}

\begin{remark}[records]\label{rem:records}
Ranks of individual curves can be large. Elkies found a curve of rank at least $28$ in 2006; by Dujella's table of records (accessed October 2026), the record
then moved to $29$ (Elkies and Klagsbrun, 2024), $30$ and $31$ (Alp\"oge and Howell, 2026) and $32$ (Elkies, Klagsbrun and Li, 2026). Whether ranks are bounded
is open. A heuristic of Park, Poonen, Voight and Wood (2019) predicts that only finitely many curves over $\Q$ have rank greater than $21$, which would
make the recent records all the more remarkable.
\end{remark}

\begin{remark}[partial Euler products]\label{rem:goldfeld}
Can the original experiment be made a theorem? Goldfeld proved in 1982 that if $\prod_{p\le x}N_p/p\sim C(\log x)^r$ for some constants $C>0$ and $r$, then
$L(E,s)$ satisfies the analogue of the Riemann hypothesis (all its non-trivial zeros on $\operatorname{Re}s=1$), $r$ is the order of vanishing of $L(E,s)$ at
$s=1$, and $C$ is determined by the leading Taylor coefficient there, with an unexpected extra factor $\sqrt2$. So the asymptotic behind
Figure~\ref{fig:bsd} is stronger than the Riemann hypothesis for $L(E,s)$, and it is not known for any curve. K.~Conrad (2005) explains where the
$\sqrt2$ comes from.
\end{remark}

\paragraph{Why it is hard.} The conjecture links two worlds with no known bridge in general. On one side is $L(E,s)$, defined by the local data $N_p$ and
continued past its region of convergence by modularity, an analytic miracle. On the other is the rank, a global object computed by descent, with $\operatorname{Sha}(E)$ as the
obstruction that descent cannot see through. The only known bridges, Heegner points (Gross--Zagier, Kolyvagin) and the Euler systems that generalise them,
produce at most one point of infinite order at a time: they work when the analytic rank is $1$ (one point to find) or $0$ (no point, but a bound). For
rank $2$ nobody knows a construction that produces two independent points from $L''(E,1)\ne0$.

\section{Reading a paper: Wiles's Millennium essay}

Andrew Wiles's \emph{The Birch and Swinnerton-Dyer conjecture}, written for the Clay Mathematics Institute's book on the Millennium Problems (2006), is the
official statement of the problem, and in ten pages it is also a history. Read it with this chapter at hand. Its first part goes from Diophantus and the
congruent number problem to Mordell's theorem; you know this from Tome~12. Then comes the definition of $L(E,s)$ and the conjecture with its full formula; check
each symbol against the box above. The last part lists the theorems: Coates--Wiles, Gross--Zagier, Kolyvagin, and Tunnell's criterion for congruent numbers.
The idea worth stealing is the shape of the conjecture itself: a local--global principle in which the local factors $p/N_p$, each trivially computable,
multiply to something that knows the global group $E(\Q)$. The original papers, B.~Birch and P.~Swinnerton-Dyer, \emph{Notes on elliptic curves I and II},
are worth opening for their tables and graphs, the first experimental evidence in this book's sense for one of the great conjectures.

\section{Open doors}

\begin{itemize}
\item Extend Figure~\ref{fig:bsd} to $x=10^8$ with a compiled program (count points with the Legendre symbol, or faster, with baby-step giant-step on the
group). Fit slopes on the last decade. How close to $3$ does the rank-$3$ curve get?
\item For the congruent number curves $E_n$, $n\le1000$ squarefree, compute $\log\prod_{p\le10^6}N_p/p$ and compare with what is known about the rank (by
Tunnell's criterion, Tome~12). Can you see the rank in the product?
\item Compute $L(E_n,1)/\Omega$ for the squarefree $n\le200$ with $n\equiv1,2,3\pmod8$ with the series of the notebook. Which values are $0$? Which non-zero values
are not of the form $\frac1{|T|^2}$ times a small integer?
\item Repeat Figure~\ref{fig:cm} for $y^2=x^3+1$ and the split primes $p\equiv1\pmod3$ (Exercise~\ref{ex:cubicD}). Which arcsine-like law appears?
\end{itemize}

\section{Exercises}

\begin{exercise}[$\star$]\label{ex:handcount}
Count the points of $y^2=x^3-x$ over $\mathbb F_5$, $\mathbb F_{13}$ and $\mathbb F_{17}$ by hand, list the points of the subgroup of order $8$ of
Lemma~\ref{lem:eightdiv} for $p=13$, and compare with Theorem~\ref{thm:countEn}.
\end{exercise}

\begin{exercise}[$\star$]\label{ex:bigp}
Use Theorem~\ref{thm:countEn} to find the number of points of $y^2=x^3-x$ and of $y^2=x^3-4x$ over $\mathbb F_{101}$ and over $\mathbb F_{1009}$
($1009=15^2+28^2$). Which of the four counts are divisible by $8$, and why?
\end{exercise}

\begin{exercise}[$\star\star$]\label{ex:jac13}
Let $p=13$ and $\chi$ the character of order $4$ with $\chi(2)=i$ ($2$ generates $\mathbb F_{13}^\times$). Tabulate $\chi$, compute $J(\chi,\rho)$, and check
Lemmas~\ref{lem:jacabs} and~\ref{lem:quarticap}.
\end{exercise}

\begin{exercise}[$\star\star$]\label{ex:quarticD}
Let $p\equiv1\pmod4$, $\chi$ of order $4$, and $D\not\equiv0\pmod p$. Show that $y^2=x^3-Dx$ has $a_p=-2\operatorname{Re}\bigl(\overline\chi(D)J(\chi,\rho)\bigr)$.
Deduce that every curve $y^2=x^3-Dx$ has $a_p\in\{\pm2a,\pm2b\}$ with $p=a^2+b^2$, and that among $D\in\{1,-1,2,-2\}$, the values $D=\pm2$ give $\pm2b$ exactly when
$p\equiv5\pmod8$.
\end{exercise}

\begin{exercise}[$\star\star$]\label{ex:cubicD}
Let $p\equiv1\pmod3$, $\psi$ a character of order $3$ of $\mathbb F_p^\times$, and $D\not\equiv0\pmod p$. Show that $y^2=x^3+D$ has
$a_p=-2\operatorname{Re}\bigl(\psi(D)\rho(D)J(\psi,\rho)\bigr)$, and that $|a_p|\le2\sqrt p$. (For $D=1$ this is Tome~12, Exercise~13.3; for $p\equiv2\pmod3$,
$a_p=0$ by Tome~12, Exercise~12.5.)
\end{exercise}

\begin{exercise}[$\star\star$]\label{ex:anbound}
For $E:y^2+y=x^3-x^2$, with $a_p$ for $p\le11$ from the table of Tome~12, Proposition~14.11 (and $a_{11}=1$), compute $a_n$ for $n\le12$ from
Proposition~\ref{prop:eulerconv}, and compare with the first coefficients of $q\prod_{n\ge1}(1-q^n)^2(1-q^{11n})^2$. Show that equality $|a_{p^k}|=(k+1)p^{k/2}$
in the bound of the proof can only happen when $a_p^2=4p$, and that this never happens for $p$ good.
\end{exercise}

\begin{exercise}[$\star\star$]\label{ex:inertmertens}
For $E_1:y^2=x^3-x$, show that $\prod_{p\le x,\,p\equiv3\,(4)}\frac{N_p}p\sim C\sqrt{\log x}$ for some $C>0$, using Tome~5, Corollary~4.5. Since $E_1$ has rank
$0$, what does the heuristic of this chapter then predict for the product over the split primes $p\equiv1\pmod4$, $p\le x$? Check it on the data of the
notebook.
\end{exercise}

\begin{exercise}[$\star\star$]\label{ex:cmmoments}
Let $t_p=a_p/\sqrt p=2\cos\theta_p$. Compute $\mathrm E[t^2]$ and $\mathrm E[t^4]$ under the law $\frac12\delta_{\pi/2}+\frac12\frac{d\theta}\pi$ of
Remark~\ref{rem:hecke} and under the Sato--Tate law $\frac2\pi\sin^2\theta\,d\theta$. Which moment separates them? Compare with the averages of $t_p^2$ and $t_p^4$
for $y^2=x^3-x$ over the odd primes below $2\cdot10^5$, which are $0.9955$ and $2.982$.
\end{exercise}

\section{Hints}

\begin{hint}{ex:handcount}
List the squares modulo $p$ first; the points of order $2$ and $O$ are always there. For $p=13$, $i=5$.
\end{hint}
\begin{hint}{ex:bigp}
$101=1^2+10^2$. For $y^2=x^3-4x$ use $n=2$ and the value of $\bigl(\frac2p\bigr)$.
\end{hint}
\begin{hint}{ex:jac13}
The powers of $2$ modulo $13$ are $2,4,8,3,6,12,11,9,5,10,7,1$.
\end{hint}
\begin{hint}{ex:quarticD}
Follow the proof of Lemma~\ref{lem:quarticap} and substitute $u=Dv$. For the second part, $\overline\chi(D)$ is a power of $i$; $\chi(-1)=-1$ when
$p\equiv5\pmod8$.
\end{hint}
\begin{hint}{ex:cubicD}
Count by $u=x^3$: the number of $x$ with $x^3=u$ is $1+\psi(u)+\overline\psi(u)$. Then substitute $u=-Dv$.
\end{hint}
\begin{hint}{ex:anbound}
$a_4=a_2^2-2$, $a_8=a_2a_4-2a_2$, $a_9=a_3^2-3$, $a_{11}=1$.
\end{hint}
\begin{hint}{ex:inertmertens}
$\log\frac{p+1}p=\frac1p+O(\frac1{p^2})$.
\end{hint}
\begin{hint}{ex:cmmoments}
$\mathrm E[(2\cos\theta)^{2k}]$ under the uniform law is the constant term of $(e^{i\theta}+e^{-i\theta})^{2k}$, a binomial coefficient.
\end{hint}

\section{Solutions}

\begin{solution}{ex:handcount}
\idea Tabulate $x^3-x$, mark the non-zero squares, and count two points for each.
\writeup Modulo $5$ the non-zero squares are $1,4$, and $x^3-x$ takes the values $0,0,1,4,0$ for $x=0,\dots,4$: three points with $y=0$, two values that are
squares, so $3+2\cdot2=7$ affine points and $N_5=8$. Modulo $13$ the non-zero squares are $1,3,4,9,10,12$, and $x^3-x$ for $x=0,\dots,12$ is
\[
0,0,6,11,8,3,2,11,10,5,2,7,0;
\]
only $3$ (at $x=5$) and $10$ (at $x=8$) are squares, so $3+2\cdot2=7$ affine points and $N_{13}=8$. Modulo $17$ the non-zero squares are
$1,2,4,8,9,13,15,16$, and $x^3-x$ is a non-zero square for $x=4,5,7,10,12,13$, so $3+2\cdot6=15$ affine points and $N_{17}=16$. Theorem~\ref{thm:countEn}
gives $5+1-(-2)=8$, $13+1-6=8$ and $17+1-2=16$. For $p=13$, $i=5$ and $P=(5,1-5)=(5,9)$; then $\langle P\rangle=\{O,(5,9),(0,0),(5,4)\}$ and the other coset
is $(1,0)+\langle P\rangle=\{(1,0),(8,7),(12,0),(8,6)\}$. That is all $8$ points: $E(\mathbb F_{13})\cong\Z/4\times\Z/2$.
\end{solution}

\begin{solution}{ex:bigp}
\idea Find the odd square, fix its sign modulo $4$, then twist.
\writeup $p=101=1^2+10^2$, $\frac{p+1}2=51\equiv3\pmod4$, so $a=-1$, $a_{101}=-2$ and $N_{101}=104=8\cdot13$ for $y^2=x^3-x$. For $y^2=x^3-4x$ ($n=2$),
$\bigl(\frac2{101}\bigr)=-1$ since $101\equiv5\pmod8$, so $a_{101}=2$ and $N_{101}=100$. For $p=1009=15^2+28^2$, $\frac{p+1}2=505\equiv1\pmod4$, and of $\pm15$
it is $-15\equiv1$, so $a_{1009}=-30$ and $N_{1009}=1040=8\cdot130$; since $1009\equiv1\pmod8$, $\bigl(\frac2{1009}\bigr)=1$ and $y^2=x^3-4x$ has $1040$
points too. The count $100$ is divisible by $4$ (the three points of order $2$ and $O$) but not by $8$: Lemma~\ref{lem:eightdiv} is about $y^2=x^3-x$, and
$y^2=x^3-4x$ is isomorphic to it over $\mathbb F_p$ only when $2$ is a square modulo $p$.
\end{solution}

\begin{solution}{ex:jac13}
\idea Read $\chi$ off the discrete logarithm.
\writeup From the powers of $2$, $\chi(u)=i^k$ when $u\equiv2^k$: $\chi=1,i,1,-1,i,i,-i,-i,1,-1,-i,-1$ at $u=1,\dots,12$, and $\rho=\chi^2$. The terms
$\chi(u)\rho(1-u)$ for $u=2,\dots,12$ ($u=0,1$ give $0$) are $i,-1,-1,i,-i,i,i,-1,-1,-i,1$, with sum $J(\chi,\rho)=-3+2i$. Then $|J|^2=9+4=13$, as
Lemma~\ref{lem:jacabs} says, and $-2\operatorname{Re}J=6=a_{13}$, as Lemma~\ref{lem:quarticap} says ($N_{13}=8$, Exercise~\ref{ex:handcount}).
\end{solution}

\begin{solution}{ex:quarticD}
\idea The substitution $u=Dv$ moves $D$ out of the sum as a character value, a fourth root of unity.
\writeup As in Lemma~\ref{lem:quarticap}, $-a_p=\sum_x\rho(x)\rho(x^2-D)=\sum_u(1+\rho(u))\chi(u)\rho(u-D)$. With $u=Dv$, $\rho(u-D)=\rho(D)\rho(v-1)=\rho(D)\rho(1-v)$
and $\chi(u)=\chi(D)\chi(v)$, so the first part is $\chi(D)\rho(D)J(\chi,\rho)=\chi^3(D)J=\overline\chi(D)J(\chi,\rho)$, and the second, with $\overline\chi$ in place of
$\chi$, is its complex conjugate. Hence $a_p=-2\operatorname{Re}\bigl(\overline\chi(D)J(\chi,\rho)\bigr)$. With $J=A+Bi$ and $\overline\chi(D)\in\{\pm1,\pm i\}$, the real
part is $\pm A$ or $\pm B$, so $a_p\in\{\pm2a,\pm2b\}$ for every $D$. For $D=\pm1,\pm2$: $\chi(-1)=\chi(g^{(p-1)/2})=i^{(p-1)/2}=(-1)^{(p-1)/4}$, and $\chi(2)^2=\rho(2)$.
If $p\equiv5\pmod8$, then $\rho(2)=-1$, so $\chi(2)=\pm i$ and $\chi(-2)=\mp i$ are imaginary while $\chi(\pm1)=\pm1$ are real: $D=\pm1$ give $\pm2a$ and $D=\pm2$
give $\pm2b$. (For $p=13=3^2+2^2$ the four values are $6,-6,-4,4$.) If $p\equiv1\pmod8$, all four values of $\chi$ are real ($\chi(2)=\pm1$), and all four
curves have $a_p=\pm2a$.
\end{solution}

\begin{solution}{ex:cubicD}
\idea Sum over the cubes instead of the squares.
\writeup $-a_p=\sum_x\rho(x^3+D)=\sum_u\bigl(1+\psi(u)+\overline\psi(u)\bigr)\rho(u+D)$, since $x^3=u$ has $1+\psi(u)+\overline\psi(u)$ solutions ($1$ for $u=0$, and $3$
or $0$ otherwise). The first part is $\sum_u\rho(u+D)=0$. In the second, put $u=-Dv$: $\psi(u)=\psi(-1)\psi(D)\psi(v)$ with $\psi(-1)=1$ (it is a cube root of unity
equal to its own cube), and $\rho(u+D)=\rho(D)\rho(1-v)$, giving $\psi(D)\rho(D)J(\psi,\rho)$; the third part is its conjugate. So
$a_p=-2\operatorname{Re}\bigl(\psi(D)\rho(D)J(\psi,\rho)\bigr)$, and since $\psi$, $\rho$ and $\psi\rho$ (of order $6$) are non-trivial, $|J(\psi,\rho)|=\sqrt p$
by Lemma~\ref{lem:jacabs}, hence $|a_p|\le2\sqrt p$. The notebook checks the formula for $p<40$ and $D=1,2,3,5$.
\end{solution}

\begin{solution}{ex:anbound}
\idea Run the recursion, and look at when the triangle inequality is an equality.
\writeup With $a_2=-2$, $a_3=-1$, $a_5=1$, $a_7=-2$, $a_{11}=1$: $a_4=a_2^2-2=2$, $a_8=a_2a_4-2a_2=0$, $a_9=a_3^2-3=-2$, $a_6=a_2a_3=2$, $a_{10}=a_2a_5=-2$,
$a_{12}=a_4a_3=-2$. So $a_1,\dots,a_{12}$ are $1,-2,-1,2,1,2,-2,0,-2,-2,1,-2$, and expanding $q\prod_{n\ge1}(1-q^n)^2(1-q^{11n})^2$ up to $q^{12}$
gives the same twelve numbers. For the bound, $|b_k|=|\sum_{j=0}^k\alpha^j\overline\alpha^{k-j}|\le(k+1)p^{k/2}$ with equality only if all the terms have the same
argument, which for $k\ge1$ forces $\alpha/\overline\alpha=1$, that is $\alpha$ real, $\alpha=\pm\sqrt p$ and $a_p=\pm2\sqrt p$. But $a_p$ is an integer and $4p$ is
not a square, so this never happens.
\end{solution}

\begin{solution}{ex:inertmertens}
\idea At an inert prime $N_p=p+1$, and Mertens in progressions does the rest.
\writeup For $p\equiv3\pmod4$, $N_p=p+1$ (Theorem~\ref{thm:countEn}), so $\log\prod\frac{N_p}p=\sum_{p\le x,\,p\equiv3\,(4)}\log\bigl(1+\frac1p\bigr)
=\sum\frac1p+c+o(1)$ with $c=\sum(\log(1+\frac1p)-\frac1p)$ convergent. By Tome~5, Corollary~4.5, $\sum_{p\le x,\,p\equiv3\,(4)}\frac1p=\frac12\log\log x+M(4,3)+o(1)$,
so the product is $\sim C\sqrt{\log x}$. Since $E_1$ has rank $0$, the heuristic says that the whole product stays bounded, so the split primes must
compensate: $\prod_{p\le x,\,p\equiv1\,(4)}\frac{N_p}p\approx\frac{C'}{\sqrt{\log x}}$. In the data, the inert product divided by $\sqrt{\log x}$ is $0.9857$,
$0.9853$, $0.9852$, $0.9852$ at $x\approx10^3,10^4,10^5,2\cdot10^5$, and the split product times $\sqrt{\log x}$ is $2.34$, $2.33$, $2.07$, $2.51$: bounded
in the range, as predicted, but fluctuating, as Figure~\ref{fig:bsd} already showed for the full product. The split part is the hard part of the problem.
\end{solution}

\begin{solution}{ex:cmmoments}
\idea Expand $(2\cos\theta)^{2k}=(e^{i\theta}+e^{-i\theta})^{2k}$.
\writeup Under the uniform law on $[0,\pi]$, the mean of $(e^{i\theta}+e^{-i\theta})^{2k}$ is its constant term $\binom{2k}k$: $2$ for $k=1$, $6$ for $k=2$. The
atom at $\frac\pi2$ contributes $0$. So the CM law gives $\mathrm E[t^2]=\frac12\cdot2=1$ and $\mathrm E[t^4]=\frac12\cdot6=3$. Under Sato--Tate, $t=2\cos\theta$ has
the semicircle law on $[-2,2]$, whose even moments are the Catalan numbers: $\mathrm E[t^2]=1$, $\mathrm E[t^4]=2$. The second moment does not separate the
laws (compare Tome~12, Exercise~13.4, where $a_p^2/p$ is $1$ on average over all curves modulo $p$); the fourth does: $3$ against $2$. The data give $0.9955$ and $2.982$, clearly the CM law.
\end{solution}

\section*{Sources and further reading}

\begin{itemize}
\item B.~J.~Birch and H.~P.~F.~Swinnerton-Dyer, \emph{Notes on elliptic curves I}, Journal f\"ur die reine und angewandte Mathematik 212 (1963), 7--25,
and \emph{Notes on elliptic curves II}, same journal, 218 (1965), 79--108.
\item A.~Wiles, \emph{The Birch and Swinnerton-Dyer conjecture}, in J.~Carlson, A.~Jaffe and A.~Wiles (eds.), \emph{The Millennium Prize Problems}, Clay
Mathematics Institute and American Mathematical Society, 2006, 31--41. This is also the source for the early history and the EDSAC~2 computations.
\item K.~Ireland and M.~Rosen, \emph{A Classical Introduction to Modern Number Theory}, 2nd edition, Springer, 1990, Chapters~8 and~18, for Jacobi sums and
point counts on $y^2=x^3-Dx$. On Gauss's diary: J.~Gray, \emph{A commentary on Gauss's mathematical diary, 1796--1814, with an English translation},
Expositiones Mathematicae 2 (1984), 97--130.
\item E.~Hecke, \emph{Eine neue Art von Zetafunktionen und ihre Beziehungen zur Verteilung der Primzahlen}, Mathematische Zeitschrift 1 (1918), 357--376,
and 6 (1920), 11--51. M.~Deuring, \emph{Die Zetafunktion einer algebraischen Kurve vom Geschlechte Eins}, Nachrichten der Akademie der Wissenschaften in
G\"ottingen, 1953. J.-P.~Serre, \emph{Quelques applications du th\'eor\`eme de densit\'e de Chebotarev}, Publications Math\'ematiques de l'IH\'ES 54 (1981),
123--201.
\item J.~Coates and A.~Wiles, \emph{On the conjecture of Birch and Swinnerton-Dyer}, Inventiones Mathematicae 39 (1977), 223--251. B.~Gross and D.~Zagier,
\emph{Heegner points and derivatives of $L$-series}, Inventiones Mathematicae 84 (1986), 225--320. V.~A.~Kolyvagin, \emph{Finiteness of $E(\Q)$ and Sha$(E,\Q)$
for a subclass of Weil curves}, Izvestiya Akademii Nauk SSSR 52 (1988). J.~Buhler, B.~Gross and D.~Zagier, \emph{On the conjecture of Birch and
Swinnerton-Dyer for an elliptic curve of rank $3$}, Mathematics of Computation 44 (1985), 473--481.
\item M.~Bhargava and A.~Shankar, \emph{The average size of the $5$-Selmer group of elliptic curves is $6$, and the average rank is less than $1$},
arXiv:1312.7859 (2013). M.~Bhargava, C.~Skinner and W.~Zhang, \emph{A majority of elliptic curves over $\Q$ satisfy the Birch and Swinnerton-Dyer conjecture},
arXiv:1407.1826 (2014).
\item A.~Dujella, \emph{History of elliptic curves rank records}, web page at the University of Zagreb (accessed October 2026). J.~Park, B.~Poonen,
J.~Voight and M.~M.~Wood, \emph{A heuristic for boundedness of ranks of elliptic curves}, Journal of the European Mathematical Society 21 (2019),
2859--2903.
\item D.~Goldfeld, \emph{Sur les produits partiels eul\'eriens attach\'es aux courbes elliptiques}, Comptes Rendus de l'Acad\'emie des Sciences de Paris
294 (1982), 471--474. K.~Conrad, \emph{Partial Euler products on the critical line}, Canadian Journal of Mathematics 57 (2005), 267--297.
\item The curves of Figure~\ref{fig:bsd} and their conductors are in J.~E.~Cremona, \emph{Algorithms for Modular Elliptic Curves}, 2nd edition, Cambridge
University Press, 1997, and in the LMFDB. All point counts, $L$-values and periods in this chapter were computed for this book.
\end{itemize}
